\chapter{CONCLUSIONS AND RECOMMENDATIONS}

\begin{thesisprose}

\section{Summary}

This study aimed to develop a portable acoustic-optical device that could classify Sweet 16 F1 watermelons as ripe or unripe before harvest. Manual assessment by tapping and visual inspection depends heavily on the experience of the harvester, while destructive methods such as cutting cannot be applied to fruit that is still to be sold. The device was designed to provide a non-destructive, on-site alternative that farmers could use directly in the field.

The study focused on the Sweet 16 F1 variety, the principal cultivar grown at the Brgy. Balubal farm in Cagayan de Oro City, and was limited to the pre-harvest assessment of one watermelon at a time while still attached to the vine. Each watermelon was classified only as ripe or unripe, with Inconclusive reported as a result rather than a third class. Classification was based on the sound of a solenoid-driven tap and a photograph of the rind, while post-harvest quality, storage behavior, and shelf life fell outside the scope. Since the system was developed and validated under the conditions of a single farm, its performance may differ for other varieties, farm settings, or regions.

The study had three objectives. The first was to design an acoustic-optical system that extracts acoustic features from the thump response and captures images of the rind with a camera. The second was to build a portable device that classifies watermelons in real time and to assess its usability. The third was to evaluate the system's ripeness classification accuracy against cut-fruit inspection and its agreement with the farmer's assessment.

Development followed the modified Waterfall model, which allowed the hardware and software to be refined as issues emerged during testing. The handheld device was built around an ESP32-S3 microcontroller with an OV5640 camera, an INMP441 microphone, a solenoid thumper, a TFT display, and two 18650 lithium-ion batteries, weighing about 0.388 kg with an estimated component cost of \pesos{}1,490. On the device, the firmware classified the rind photograph and the tap recording separately, then combined the two predictions through log-likelihood ratio weighting and displayed the individual classifications alongside the final verdict. A verdict of Inconclusive indicated that another assessment of the same fruit was needed.

The acoustic dataset consisted of 494 tap recordings, with ripe samples collected at Cogon Public Market and unripe samples at a farm in Brgy. Balubal, while the image dataset consisted of 485 rind photographs taken with smartphones. For the acoustic branch, the thump response was isolated with an adaptive window, and 31 spectral, time-domain, and cepstral features were extracted from it. For the image branch, the photographs were cleaned through cropping, glare reduction, sharpening, and lighting and color normalization, then resized to 96 \ensuremath{\times} 96 pixels.

On the acoustic side of the first objective, seven classifiers were trained using 395 original tap recordings and two noise-augmented copies of each, for a total of 1,185 training samples. Tuning the Decision Tree increased its mean ripe-class F1-score across five validation folds from 0.911 to 0.954 and raised its F1-score on the 99 held-out test recordings from 0.933 to 0.984, with an accuracy of 0.990 and a ripe precision of 1.000. The tuned Decision Tree was selected for deployment, since it tied with the default XGBoost for the highest test F1-score while fitting on the ESP32-S3 as a 1.66 KB header file. Noise augmentation also improved the average F1-score on the noisy test sets for six of the seven classifiers.

On the image side, eight architectures were each trained with 17 optimizer configurations, for a total of 136 runs, on 485 images divided into 388 training, 48 validation, and 49 test images, with the run for each architecture selected by its validation loss. EfficientNet-B0 correctly classified 47 of the 49 held-out images (95.92\%), including all 24 unripe images, and had the highest ROC AUC among the selected configurations at 0.9967. After conversion to INT8 for deployment on the ESP32-S3, it correctly classified 48 of the 49 images, for an accuracy of 97.96\% and a macro F1-score of 0.9796, while EfficientNet-Lite0, which had tied it before conversion, fell to 40 correct images.

The second objective produced a handheld device that guided the user through image capture and tapping and performed all classification on the ESP32-S3, without an internet connection or external computer. Although it was designed to run on two 18650 batteries, it was powered by a power bank during field testing due to hardware limitations. In the usability evaluation with 20 farmers and harvesters, it received a mean SUS score of 75.63, which falls within the Good (Grade B) range. Its average processing time per watermelon was 16,944.82 ms (about 16.9 seconds), with 16,920.90 ms for the image model together with the preprocessing and 23.86 ms for the acoustic model.

Field testing for the third objective covered 43 Sweet 16 watermelons at the Brgy. Balubal farm, with 46 assessments in total. The system gave a definitive verdict on the first attempt for 40 of them, and the three Inconclusive results were all resolved on retrial. Against destructive ground truth, it correctly classified 41 of the 43 watermelons (95.35\%) with a ripe precision of 100\%, compared with 97.67\% for the farmer's manual assessment, and its agreement with the farmer reached a Cohen's kappa of 0.780. Counting the Inconclusive results as errors, the initial error rate was 13.04\% and the re-trial frequency was 6.52\%.

\section{Conclusion}

The first objective of designing the acoustic-optical system was accomplished. The device captured the thump response through the INMP441 microphone and the rind through the OV5640 camera, and extracted 31 acoustic features on the ESP32-S3 using filterbank and DCT parameters exported from Python. The selected Decision Tree and EfficientNet-B0 models both performed strongly on held-out data, and noise augmentation improved the robustness of most acoustic classifiers to background noise. However, the Decision Tree relied mainly on spectral centroid, and since the ripe and unripe recordings were collected at different sites and under different conditions, part of this separation may reflect the collection conditions rather than ripeness alone. In addition, the image and audio data were not paired by fruit, individual fruits were not tracked, and the features calculated on the device were not checked against those calculated in Python for the same recordings.

The second objective of developing a portable real-time device and assessing its usability was also accomplished. The device guided the user through image capture and tapping and displayed the verdict on the device itself, without an internet connection or external computer, although it was powered by a power bank rather than its own batteries during field testing. Its mean SUS score of 75.63 falls within the Good (Grade B) range, indicating acceptable usability, and its average processing time of about 16.9 seconds per watermelon allows each fruit to be assessed on the spot, without sending data to an external computer. Most respondents found it easy to use without technical support or prior learning, with 19 of 20 disagreeing that they would need a technical person's help and all 20 disagreeing that they needed to learn a lot before getting started.

The third objective of evaluating the system's classification accuracy and its agreement with the farmer was also accomplished. The system classified watermelons on the vine with an accuracy of 95.35\% against destructive ground truth, approaching the farmer's 97.67\%, and like the farmer, it did not classify any unripe watermelon as ripe. Its substantial agreement with the farmer (\ensuremath{\kappa} = 0.780) also indicates that its verdicts are consistent with current harvesting practice. These results, however, come from a single farm, a single variety, and only 8 unripe watermelons, so a single misclassification changes the unripe results considerably. The field recordings were also quieter than the training recordings, and the fusion rates were preset rather than measured from data, both of which should be addressed before wider deployment.

Overall, the study shows that combining acoustic and optical analysis on a low-cost microcontroller can provide a portable, non-destructive aid for judging watermelon ripeness before harvest. While the farmer's experience remains slightly more accurate, the system reached comparable accuracy and did not classify any unripe fruit as ripe, which makes it a promising decision-support tool for growers. If the limitations above are addressed, the device has the potential to support harvest decisions on a wider scale, particularly for growers who lack years of experience in judging ripeness by hand.

\section{Recommendation}

Based on the findings and limitations observed during the design, development, and evaluation of the portable acoustic-optical watermelon ripeness detection system, the following recommendations are proposed to improve its classification accuracy and hardware durability, extend its use to a wider range of varieties, farms, and growing conditions, and support its adoption as a practical pre-harvest tool for watermelon growers.

The device should be made moisture-resistant, since field use exposes it to dew, rain, and wet rinds. A sealed enclosure with protected openings for the microphone and camera, along with water-resistant connectors and buttons, would protect the electronics without blocking the sensors. The device could also be made smaller and lighter through a custom printed circuit board and a more compact battery and thumper arrangement, making it easier to hold in one hand and carry through the field. The power system should also be revised so that the device can run on its own batteries during field use, as it relied on a power bank during field testing, and its battery life over a full day of harvesting should be measured.

The user interface should be simplified, since some respondents were neutral about the consistency of its results and about using it frequently. Clear icons, color-coded verdicts, and instructions in the local language would make the device easier to learn, while sound or vibration feedback could confirm each step without the user having to read the screen. A brighter or anti-glare display would also improve readability under direct sunlight. Guidance on how firmly to press the barrel against the fruit would help keep each tap consistent between users.

Future data collection should record both classes at the same site, in the same fruit position, and under the same conditions, so that the classifiers learn ripeness rather than differences between recording environments. Each watermelon should also be tracked individually, with its image and tap recordings collected together, so that the true and false positive rates used in the fusion can be measured from paired data instead of being preset. Data should also be collected across different growing seasons, farms, watermelon varieties, times of day, and weather conditions, so that the system remains usable outside the season and conditions covered in this study.

Background noise from more field environments should be recorded for augmentation, and the end of the thump window could be detected with a rule that is less sensitive to noise. The features computed on the ESP32-S3 should also be compared with those from librosa on identical recordings, and a Decision Tree that outputs a confidence score could be exported so that weak acoustic predictions carry less weight in the fusion.

For the image branch, the model should be retrained or fine-tuned on images captured by the device's own camera, since the training images in this study were taken with smartphones. The images should also cover a wider range of lighting, including early morning, midday, and overcast conditions, so that the model is less affected by changes in daylight. The processing time of the image branch, which accounted for nearly all of the 16.9 seconds per watermelon, could also be reduced by using a smaller input size, a lighter architecture, or fewer preprocessing steps.

The system's capabilities could also be extended beyond binary classification. Future versions could distinguish more ripeness stages, such as unripe, nearly ripe, ripe, and overripe, or estimate sugar content against Brix readings from a refractometer. A companion mobile application connected through Bluetooth could let farmers keep a record of assessed fruits and track ripening across their field.

Future development of the system could adopt an Agile methodology instead of the modified Waterfall model. Since the hardware, firmware, and models in this study required repeated revision as field issues emerged, short development cycles with regular field testing would allow problems to be identified and corrected earlier.

Beyond further development, the device could be brought closer to the farmers it was designed for through technology transfer. Partnering with the Department of Agriculture, local government units, or farmer cooperatives would allow the device to be piloted on more farms, and short orientation sessions with simple instructions in the local language would help new users get started, even though most respondents did not feel they needed technical assistance. The cost of the device should also be kept low as it is refined, so that small-scale farmers can afford it.

\end{thesisprose}